%% ------------- English version ---------------
\documentclass[english,hidelinks]{sbrt}
\usepackage{newtxtext}
\usepackage[english]{babel}
\usepackage[utf8]{inputenc}
\newtheorem{theorem}{Theorem}
%% ---------------------------------------------
\setlength{\marginparwidth}{2cm}
\usepackage{todonotes}
\newcommand\rigel[1]{\todo[color=green, inline]{\textbf{Rigel}: #1}}
\usepackage{orcidlink}
\usepackage{hyperref}
\usepackage[
  compatibility=false,
  font=footnotesize,
  labelsep=period,
  format=plain,
  justification=raggedright,
  singlelinecheck=false
]{caption}

\begin{document}

\title{Low-Cost Embedded System for Beehive Monitoring}

\author{Igor Costa, Jorge Alves, Ian Dias and Davi Ito
  \thanks{The authors are with Ibmec, Rio de Janeiro, RJ, Brazil.
  E-mails: igorcosta.silva10a@gmail.com,
  jorgedocruzeiro07@gmail.com,
  iandpcarvalho@gmail.com,
  cdavi8249@gmail.com.}%
  }

\maketitle

\markboth{XLIV BRAZILIAN SYMPOSIUM ON TELECOMMUNICATIONS AND SIGNAL PROCESSING - SBrT 2026, SEPTEMBER 29TH TO OCTOBER 2ND, 2026, SALVADOR, BA}{}

\begin{abstract}
\end{abstract}
\begin{keywords}
Beehive monitoring, embedded systems, Internet of Things, computer vision.
\end{keywords}

\section{Introduction}

Bees play a central role in natural and agricultural ecosystems because pollination by bees supports the reproduction of flowering plants and the production of crops. The western honey bee (\textit{Apis mellifera}) is among the most widely distributed and frequently observed crop pollinators, while its importance in natural habitats can vary according to the ecological context. At the same time, bee populations face multiple pressures, including habitat loss, invasive species, emerging diseases, pesticide exposure, and climate change. These conditions strengthen the need for reliable observation of colony and bee activity, especially when changes must be identified over time rather than through isolated inspections \cite{hung2018pollinators,brown2009conservation}.

Traditional hive inspections provide valuable information about colony conditions, but they are intermittent, require human presence, and may disturb the colony when the hive is opened repeatedly. Embedded systems are computing platforms designed for specific functions and commonly optimized for constraints such as processing time, reliability, power consumption, size, and cost. Internet of Things (IoT) architectures extend these systems by connecting sensors and other physical devices to communication networks and data services, allowing data to be processed locally, remotely, or across both levels. De Micco et al. reviewed the characteristics and applications of embedded systems, including their relation to IoT and edge computing, while Ray surveyed IoT architectures composed of devices, sensors, communication protocols, services, and processing layers. This technological foundation motivates a low-cost, non-invasive beehive monitoring system in which environmental and visual measurements are collected by embedded devices and made available for subsequent analysis \cite{demicco2020embedded,ray2018iotarchitectures}.

\textit{Related work.} Entrance-focused computer vision has been combined with environmental sensing to measure colony activity without opening the hive. Zheng et al. placed temperature and humidity sensors at the hive center, margin, entrance, and surrounding environment and pointed a webcam at the entrance. Their system, which combines YOLOv5 with DeepSORT, reached 83.5\% $\pm$ 0.7 MOTA and 77.3\% $\pm$ 0.2 MOTP at up to 16 frames per second. Narcia-Macias et al. reported 96.28\% tracking accuracy and an F1-score of 0.8319 for pollen detection with the low-cost, YOLOv7-tiny-based IntelliBeeHive, but their mite data were insufficient and partly placeholders, so their Varroa detection is a proof of concept rather than a validated field capability \cite{zheng2024intelligent,narciamacias2024intellibeehive}. Henry et al. built a wireless sensor network that recorded temperature, relative humidity, and sound inside hives in an apiary. It shows continuous in-hive measurement working in the field, and it depends on a network that many apiaries lack \cite{henry2019precision}.

Tu et al. counted bees at the entrance with background subtraction on a Raspberry Pi and reported $R^2$ values of 0.987 for the count, 0.953 for incoming activity, and 0.888 for outgoing activity against manual annotations. Their results support counting on a low-cost platform and checking the count against human labels, although illumination, camera placement, and overlapping bees affect the result \cite{tu2016automatic}. Sledevi\v{c} and Matuzevi\v{c}ius published 7200 frames from eight hives, filmed with a smartphone 30~cm above the landing boards at 1920 by 1080 pixels and 50 frames per second, with every bee labelled by hand \cite{sledevic2024dataset}. On this material, Sledevi\v{c} et al. reached 98\% mean precision in bee detection with YOLOv8m at up to 36 frames per second and recognized foraging, fanning, washboarding, and defence with 87\% accuracy \cite{sledevic2025behavior}. Both studies use smartphone footage, and neither measures how a camera of the class found in low-cost embedded stations affects the detection.

This work proposes a low-cost station that records the entrance activity and the temperature of a beehive without opening it and without a network. A camera above the entrance captures short image windows on a fixed schedule, a sealed sensor reads the temperature inside the hive, another reads temperature and humidity outside, and both streams reach a memory card stamped by the same clock. Offline, a detector counts the bees visible in each window, and each count is stored next to the temperature of the same time, so activity can be read against the conditions that produced it. The group has no access to a hive. The detector is therefore evaluated on the public videos of Sledevi\v{c} and Matuzevi\v{c}ius, and the main experiment films those videos with the station camera to measure how much the camera degrades the count against the same labels. The station counts visible bees. The camera is off between windows, so it does not follow a bee from one window to the next, and it claims nothing about entries and exits, pollen, \textit{Varroa}, the queen, or swarming. The implementations of this work are available in \cite{repo}. Section~\ref{sec:system-description} describes the station, Section~\ref{sec:method} the method, Section~\ref{sec:results} the partial results, and Section~\ref{sec:conclusion} concludes.

\section{System Description}
\label{sec:system-description}

The station is a single ESP32-WROVER-DEV board that drives the camera, reads the sensors, and writes everything to a memory card.

\subsection{Image Acquisition}

An OV2640 camera records the entrance at 1280 by 720 pixels. The board carries external PSRAM, which the camera needs to hold a frame at that resolution. The camera sits above and slightly in front of the entrance, aimed at the landing board, and blocks neither the entrance nor the ventilation. Zheng et al. and Narcia-Macias et al. point a camera at the entrance for the same reason, and Tu et al. showed that a low-cost board is enough to count bees there \cite{zheng2024intelligent,narciamacias2024intellibeehive,tu2016automatic}. The station stores JPEG frames instead of H.264 video, so every frame keeps its own timestamp and the originals remain available for a second analysis.

\subsection{Environmental Sensing}

A DS18B20 measures the temperature inside the hive and an AHT20 measures temperature and humidity outside, in the shade next to the entrance. Both are read every three minutes. The interior sensor measures only temperature, so heat-shrink tubing can seal it completely, with no exposed contact for the propolis that bees spread over objects left in the hive. A humidity sensor needs contact with the air, so humidity is measured outside only. A three-wire cable carries power and data from the DS18B20 to the board mounted outside. Zheng et al. placed sensors inside and around the hive, and Henry et al. kept continuous temperature records inside hives in an apiary \cite{zheng2024intelligent,henry2019precision}.

\subsection{Storage, Timekeeping, and Power}

The camera, the PSRAM, and the flash leave seven free pins on the board. The microSD card takes three of them in one-bit mode, a DS3231 real-time clock and the AHT20 share the two lines of an I2C bus, and the DS18B20 uses one pin on its 1-Wire bus. The clock keeps timestamps consistent without any network. The card is retrieved about once a week and analysed on a separate computer, so a failing link cannot interrupt data collection, and remote transmission as in Henry et al. remains an extension \cite{henry2019precision}. The prototype runs from a power bank. A solar panel with a battery is left for a later version, so this work makes no claim of autonomous operation.

\section{The Method}
\label{sec:method}

\rigel{AC Semana 7:
Nesse momento já teremos a seção Introdução, a seção II e a Seção III que trata do método que o grupo propõe.
Atividade relativa a parte escrita: Elaborar a seção III que deverá tratar sobre o método proposto:

(1) Apresente, em detalhes, o funcionamento do sistema desenvolvido. Como nessa seção estão as contribuições do grupo, não é necessário inserir referências que dêem suporte;

(2) Elabore um diagrama que facilite a compreensão do método (use o LucidChart por exemplo, exporte o diagrama como PDF para manter a boa resolução no artigo, cuidado com os espaços em branco que podem sobrar às margens do diagrama (no momento de exportar como PDF faça um recorte personalizado para tirar as margens em branco indesejáveis).

De forma alternativa (ou adicionalmente), o grupo poderá gerar um pseudocódigo para mostrar como o sistema funciona (passos principais do método proposto);

(3) Um parágrafo deverá ser escrito e deverá referenciar o diagrama criado na seção III Método. Explique o diagrama criado ao leitor.
(4) Outro parágrafo deverá ser escrito e referenciar uma foto do protótipo. Explique o que está na foto ao leitor.}

\section{Results}
\label{sec:results}

\rigel{AC Semana 8:
Nesse momento já teremos a seção Introdução, a seção II e a Seção III que trata do método que o grupo propõe.
Atividade relativa a parte escrita: Elaborar a seção IV que deverá tratar sobre os resultados experimentais:

(1) Apresente, em detalhes, os resultados qualitativos e quantitativos obtidos até o momento. No mínimo resultados qualitativos.

(2) Elabore o Abstract e Conclusions.}

\rigel{AC Semana 9: Continuar na elaboração da seção IV que deverá tratar sobre os resultados experimentais:

(1) Introduza na Seção IV (Results), em detalhes, os resultados quantitativos obtidos (além de resultados qualitativos já conseguidos).

(2) Refine o abstract. Na última linha do abstract devemos inserir algum resultado que seja interessante para o leitor e que represente algum aspecto considerado como a maior contribuição do grupo

(3) Refine a seção Conclusions. Insira resultados que resumam a avaliação do sistema embarcado proposto.}

\section{Conclusion}
\label{sec:conclusion}

\bibliographystyle{ieeetr}
\bibliography{refs}

\end{document}